problem:  
Let \((M^{4n},g)\) be a compact positive quaternion‑Kähler manifold with \(\mathrm{Hol}(g)\subset \mathrm{Sp}(n)\mathrm{Sp}(1)\), \(n\ge2\), and let \(Z\) denote its twistor space. Then \(Z\) is a complex contact Fano manifold of complex dimension \(2n+1\), equipped with a holomorphic contact sequence  
\[
0 \longrightarrow \mathcal D \longrightarrow TZ \longrightarrow L \longrightarrow 0,
\quad\text{with } -K_Z = (n+1)L.
\]  
For a homogeneous (Wolf) space, the topological intersection numbers  
\[
\Theta_W(r):=\int_Z c_1(L)^{2n+1-r}\,c_r(\mathcal D)
\]
can be computed explicitly from representation theory of the corresponding compact Lie group, and they exhibit a rigid pattern depending only on the Lie type.

**Problem:**  
For a general compact positive quaternion‑Kähler manifold \((M^{4n},g)\), define
\[
\Theta_M(r):=\int_Z c_1(L)^{2n+1-r}\,c_r(\mathcal D),
\]
for \(0\le r\le 2n+1\).  
Is it true that the normalized ratios
\[
R_M(r,s):=\frac{\Theta_M(r)}{\Theta_M(s)}, \qquad (0\le r,s\le 2n+1)
\]
coincide with those of some Wolf space if and only if \(M\) is locally symmetric?  
Equivalently, do the intersection numbers of the contact distribution on the twistor space provide topological rigidity strong enough to force local symmetry?

Why is it a "good" problem:  
This question refines the search for curvature rigidity by isolating precise **intersection‑theoretic invariants** of the contact structure on the twistor space. It formulates a new, purely topological/algebraic condition—equality of integral Chern number ratios—that could characterize symmetry within the family of positive quaternion‑Kähler manifolds.  

It reaches beyond the known Picard‑number classifications and beyond the automatic Kähler–Einstein proportionalities, engaging characteristic classes that reflect the full nonlinearity of the contact structure. These intersection numbers mix geometry (through the quaternionic twistor fibration), topology (via Chern classes), and representation theory (through known values for adjoint varieties).  

A positive result would give an entirely cohomological characterization of Wolf spaces within the quaternion‑Kähler universe, bridging algebraic topology, complex contact geometry, and Riemannian holonomy. Even partial progress would develop new tools connecting the algebraic topology of contact Fano manifolds to differential geometric rigidity phenomena. It epitomizes the “narrow entrance, profound depth” ideal—straightforward to state, yet linking multiple frontiers of modern geometry.